\hypertarget{appendix-f11d1a521253}{}
\section*{Floating-point environment and compiler modes}
\section{Floating-point environment and compiler modes}\label{floating-point-environment-and-compiler-modes}

\textbf{Curriculum depth:} Developed · \textbf{First seen:} C004 · \textbf{Latest development:} C005

\subsection{Definition and mechanism}\label{definition-and-mechanism}

The environment includes rounding direction and exception status. Runtime mode changes can fail and must restore prior state. Strict compiler settings are needed; fast-math may discard required special-value or environment assumptions.

M005 uses strict arithmetic and disables implicit FMA contraction to distinguish fused from separate graphs. The full floating environment is restored after the lab; cross-platform bitwise reproducibility is not established.

\subsection{Prerequisites}\label{prerequisites}

\begin{itemize}
\tightlist
\item
  \hyperref[appendix-637d164ae5ca]{Directed rounding modes} --- The signed rounding direction is state that a demonstration or enclosure algorithm must control.
\item
  \hyperref[appendix-d6607cc4c235]{NaNs and payloads} --- Unordered comparisons and payload distinctions require explicit metric, comparison and serialization policy.
\end{itemize}

\subsection{Used by}\label{used-by}

\begin{itemize}
\item
  \hyperref[appendix-5f1a2c3f97ff]{Numerical reproducibility} --- Rounding, contraction, subnormal and optimization assumptions must match.
\item
  \hyperref[appendix-e979eb012782]{Flush-to-zero and denormals-are-zero} --- uses this mechanism in its explanation.
\item
  \hyperref[appendix-1823a14330cb]{Interval arithmetic} --- uses this mechanism in its explanation.
\end{itemize}

\subsection{Recall and next depth}\label{recall-and-next-depth}

Reconstruct the concrete example in the definition and identify its validity assumptions.

\textbf{Pending:} Extend this mechanism to later algorithms and target-specific behavior. See the M005 chapter for its current examples and validity assumptions.

\subsection{Source and navigation}\label{source-and-navigation}

\href{Continuous\%20Notes/Computing\%20Foundations\%20-\%20From\%20Bits\%20to\%20C++\%20and\%20Systems\%20APIs\#When\%20arithmetic\%20meets\%20uncertainty}{M005 development}

\href{Continuous\%20Notes/Computing\%20Foundations\%20-\%20From\%20Bits\%20to\%20C++\%20and\%20Systems\%20APIs\#A\%20stretching\%20ruler\%20for\%20real-valued\%20quantities}{M004 teaching chapter} · \hyperref[appendix-fd090d982a15]{Per-part definitions} · \href{}{Dependency Map (02-Foundations-II, appendix-4ad47bdd269a)}
